
\chapter{Introdução}

A área de compiladores tem evoluído e enfrentado desafios contínuos, com novas arquiteturas e tecnologias sendo utilizadas para executar programas cada vez maiores e mais complexos. Além disso, há uma emergente quantidade de linguagens de programação sendo criadas, e cada nova linguagem necessita de sua própria ferramenta de compilação. A complexidade envolvida na construção de tais ferramentas tem direcionado seu \emph{design} cada vez mais para modelos modulares e reutilizáveis \cite{muchnick1997advanced}. Uma divisão frequentemente vista na construção de compiladores é entre o \emph{frontend} e o \emph{backend}. O \emph{frontend} é responsável por fazer o tratamento dos detalhes específicos da linguagem fonte, enquanto no \emph{backend} são tratados os detalhes da arquitetura alvo em específico \cite{dragao}. Essa tendência em modelos modulares evidência a importância na interface desses módulos, as chamadas representações intermediárias (IRs - do inglês \emph{Intermediate Representations}).

Segundo \citeonline{dragao}, o processo de traduzir código da linguagem fonte para linguagem de máquina envolve o uso de diversas representações intermediárias. Essas representações são construídas durante o processo de compilação de um programa, funcionando tanto como estruturas que contêm informações necessárias para análises do código, quanto como interfaces de comunicação entre diferentes etapas na \emph{pipeline} de compilação. As representações intermediárias podem ser estruturas de dados armazenados em memória durante a execução do compilador, ou também representações textuais ou binárias escritas em arquivos. Dentre estas é possível citar alguns exemplos comumente encontrados na literatura, como: Árvore Sintática Abstrata, Grafo de Fluxo de Controle, Código de Três Endereços, entre outras.

Uma das principais IRs, que é utilizada durante etapas de otimização de fluxo de dados nos principais compiladores modernos de linguagens imperativas, é a forma de atribuição única estática (SSA - \emph{Single Static-Assignment}). Essa representação é uma forma que impõe que todas as atribuições em um procedimento devem ser feitas a somente uma variável, limitando que uma nova variável seja criada a cada nova atribuição, ou seja, na qual as variáveis são imutáveis. SSA também define a notação conceitual da função $\phi$, que é capaz de fazer a união de diferentes definições de variáveis no grafo de fluxo de controle. Segundo \citeonline{muchnick1997advanced}, as características da forma SSA promovem a separação entre os valores operados em um programa, dos locais nos quais estes são armazenados, tornando possíveis versões mais eficazes de diversas otimizações, principalmente de fluxo de controle e dados.

Dadas as vantagens supracitadas, SSA foi utilizada como base para a criação da representação intermediária de um dos conjuntos de ferramentas de compilação de código que está no estado da arte atualmente, o \emph{framework} LLVM (do inglês, Máquina Virtual de Baixo Nível). Esse conjunto de ferramentas foi projetado e desenvolvido com a proposta de oferecer um conjunto de ferramentas para auxiliar no desenvolvimento de novos compiladores de forma eficiente, modular e reutilizável \cite{lattner2004llvm}. Desde seu início, no ano de 2000, o uso da LLVM vem crescendo, sendo usado em uma variada gama de projetos de compiladores comerciais e pesquisas acadêmicas. A LLVM apresenta dois pontos principais: a definição de sua representação intermediária de código (LLVM-IR) que é independente da linguagem fonte, carrega informação de tipagem, controles de fluxo de controle explícitos, e que oferece uma ampla gama de funcionalidades, servindo como uma forma unificada de representar o código para análise, modificação e distribuição; e o design do compilador, que, através de seu modelo modular, utiliza a representação intermediária para fornecer uma vasta combinação de capacidades de compilação e otimização de código arbitrário \cite{lattner2004llvm}. LLVM é um projeto de código fonte aberto e que contém licenças não restritivas\footnote{\url{https://github.com/llvm/llvm-project}}, fato que contribui com sua popularidade, evolução e uso por parte da comunidade de desenvolvimento de software.

O paradigma funcional recebe essa denominação porque se baseia em programas compostos inteiramente por funções \cite{hughes1989functional}. A programação e computação feita neste paradigma é feita por meio da definição, aplicação e retorno de funções. Funções que apenas recebem argumentos e dão um retorno, sem mais efeitos além disso. Ou seja, em linguagens puramente funcionais não existem os chamados efeitos colaterais, que são as alterações de estados de memória de forma que essas alterações possam influenciar em uma computação futura. A ausência de efeitos colaterais é garantida pela imutabilidade das variáveis. E confere, dessa forma, a chamada transparência referencial, ou pureza, que tem influência profunda em como os programas são construídos e entendidos, o que leva a uma maior garantia de correção do código \cite{whymattered}. A programação através do paradigma funcional tem seus pontos fortes vastamente defendidos na literatura, por exemplo, na área da programação paralela e distribuída. A ausência de efeitos colaterais no paradigma funcional favorece a programação paralela pois, sem efeitos colaterais, problemas clássicos da programação paralela, como condição de corrida, são evitados \cite{hammond2011parallel} \cite{whymattered}.

Todavia, há efeitos colaterais que são necessários ao programar computadores. Exemplos clássicos são: a impressão na saída padrão, escrita em arquivos, inserções em banco de dados, entre outros; que linguagens puramente funcionais devem suportar. Portanto, as linguagens puramente funcionais implementam meios de simular os efeitos colaterais de forma pura: classicamente através de mônadas \cite{peyton1993imperative}, e mais recentemente por meio do uso de efeitos algébricos \cite{leijen2016algebraic}.

Compiladores de linguagens funcionais também fazem uso de diversas representações intermediárias durante o processo de compilação. Uma representação muito importante em etapas de otimização e análise de código funcional é a ANF, ou Forma Administrativa Normal (\emph{Adminstrative-Normal Form}) \cite{sabry1992reasoning}. Semelhante ao SSA, ANF tem sua importância ao deixar explícito o fluxo de controle no código e, segundo \citeonline{ssaToAnf}, tem inclusive vantagens quanto à análise do fluxo de controle quando comparado com SSA. A ANF é fundamentalmente uma restrição a termos lambda, ou seja, é expresso dentro do paradigma de programação funcional.

No que tange este trabalho, \citeonline{appel1998ssa} demonstrou a correspondência entre o paradigma funcional e a forma de atribuição única estática. Nota-se que essa correspondência pode advir do fato de que ambas não permitem atribuição a variáveis uma vez iniciadas. Além disso, \citeonline{ssaToAnf} demonstram que há um algoritmo que estaticamente faz a tradução da forma SSA para ANF. \citeonline{rigon2020inferring} exploram a correspondência entre SSA e o paradigma funcional para para inferir efeitos algébricos em uma representação funcional gerada a partir de SSA. Portanto, usando do fato da LLVM ter sua representação intermediária baseada na forma SSA há a possibilidade de explorar a tradução da LLVM-IR para o paradigma funcional em ANF.

\section{Objetivo Geral}

O objetivo do presente trabalho é investigar a possibilidade de traduzir a representação intermediária SSA gerada pela LLVM para representação funcional em ANF, tendo como objetivo futuro o uso dessa representação em uma extensão do método proposto por \citeonline{rigon2020inferring} para inferir efeitos algébricos em código imperativo real.

\section{Objetivos Específicos}

\begin{itemize}
    \item Estudar os conceitos de compiladores de tradução de código, representação intermediária, SSA e o paradigma de programação funcional;
    \item Estudar sobre a tradução de código intermediário (SSA) da LLVM para código puramente funcional;
    \item Implementar um tradutor de código SSA gerado pela LLVM para código puramente funcional;
    \item Elucidar a possibilidade de fazer tal tradução e quais são as ressalvas quanto a essa abordagem.
\end{itemize}

\section{Estrutura do Trabalho}

O presente trabalho é estruturado em mais 3 capítulos: no Capítulo \ref{cap:fundamentacao}, são explorados em profundidade os principais conceitos teóricos necessários para a proposta e implementação deste trabalho; o Capítulo \ref{cap:desenvolvimento} detalha a proposta, explica a implementação do método, apresenta os resultados e faz uma breve discussão; por fim, no Capítulo \ref{cap:conclusao}, é feita a conclusão sobre o resultado obtido neste trabalho.
